La integraci¢n efectiva de las Tecnolog¡as de la Informaci¢n y la Comunicaci¢n (TIC) en los modelos productivos del sector primario constituye uno de los desaf¡os sociot‚cnicos m s apremiantes para la competitividad agr¡cola en Colombia. En este escenario, la brecha digital trasciende la mera conectividad; se manifiesta en la ausencia de ecosistemas digitales que permitan a los peque¤os productores mitigar la asimetr¡a de informaci¢n y participar activamente en cadenas de valor globales. Bajo esta premisa, el presente art¡culo documenta el ciclo de vida de ingenier¡a y la evoluci¢n arquitect¢nica del Portal Agro-Comercial del Huila, una plataforma web dise¤ada para visibilizar la oferta productiva del municipio de Teruel y reducir la intermediaci¢n comercial.

En su fase constructiva inicial, el sistema se desarroll¢ bajo un enfoque pragm tico, priorizando el cumplimiento de requisitos funcionales mediante una arquitectura de despliegue monol¡tico. Si bien esta estrategia permiti¢ validar la propuesta de valor mediante la gesti¢n de identidades digitales y cat logos agr¡colas, la operaci¢n en entornos reales revel¢ limitaciones estructurales cr¡ticas. El fuerte acoplamiento entre la infraestructura y la l¢gica de negocio deriv¢ en ventanas de mantenimiento obligatorias y un alto riesgo de regresi¢n, factores que erosionan la confianza del usuario final en contextos donde la disponibilidad del servicio es determinante para la adopci¢n tecnol¢gica.

Ante la evidencia de estas fricciones operativas, y trascendiendo el alcance de un reporte t‚cnico convencional, este trabajo propone una refactorizaci¢n arquitect¢nica fundamentada en la inyecci¢n de \textit{Feature Toggles} (banderas de funcionalidad). La metodolog¡a adoptada articula un enfoque mixto: por un lado, el an lisis del caso de estudio real, y por otro, una revisi¢n sistem tica de la literatura reciente (2016--2025). Esta dualidad permite contrastar la rigidez del sistema legado con un modelo evolutivo que integra gesti¢n din mica de configuraci¢n en tiempo de ejecuci¢n.

El aporte central de esta investigaci¢n radica en la formalizaci¢n de un marco de trabajo que desacopla conceptual y t‚cnicamente el despliegue (la instalaci¢n de binarios) de la liberaci¢n (la activaci¢n de valor para el negocio). Se demuestra te¢ricamente c¢mo esta separaci¢n habilita la transici¢n hacia pr cticas de Entrega Continua (\textit{Continuous Delivery}), facilitando la implementaci¢n de estrategias avanzadas como lanzamientos graduales (\textit{Canary Releases}) y pruebas A/B en producci¢n. Finalmente, se aborda la gobernanza de la deuda t‚cnica, concluyendo que la sostenibilidad a largo plazo de las herramientas digitales en el agro depende de arquitecturas resilientes, capaces de evolucionar org nicamente sin interrumpir la operaci¢n cr¡tica de los agricultores.
